\section{Ablation Studies}
\label{sec:appendix_ablations}

\vspace{-0.25em}
\noindent
This appendix isolates \textbf{which components actually matter} for \textsc{CITA} and \textbf{why}.
Our central hypothesis is that \textsc{CITA} succeeds at \textbf{instruction-conditioned switching} only when the optimization remains within a \textbf{stable local neighborhood} of a reference policy while still learning \textbf{sharp, instruction-dependent preference gaps}.
Accordingly, we ablate the two parameters that directly control this trade-off:
\textbf{(i) the KL anchor weight} $\lambda_{\text{KL}}$, which governs \textbf{stability and regime co-existence}, and
\textbf{(ii) the preference temperature} $\beta$, which governs \textbf{gap sharpness and switching sensitivity}.
We report three outcome signals throughout:
\textbf{reward margin} (training separation), \textbf{ECLIPTICA score} (switchability under controlled instruction flips), and \textbf{stability} (run-to-run and within-run behavior).

% -----------------------------------------------------------------------------
\subsection{Effect of KL Weight ($\lambda_{\text{KL}}$)}
% -----------------------------------------------------------------------------

\vspace{-0.25em}
\noindent
\textbf{Why this ablation matters.}
In \textsc{CITA}, the KL term is not a cosmetic regularizer: it is the mechanism that keeps \textbf{multiple instruction regimes co-located} so that switching remains \textbf{nearby and traversable} rather than collapsing into a single dominant posture.
Too small $\lambda_{\text{KL}}$ permits \textbf{runaway preference margins} that overwrite neighboring regimes; too large $\lambda_{\text{KL}}$ imposes a \textbf{hard trust region} that prevents learning the instruction-conditioned separation needed for switching.

\begin{table*}[ht!]
\centering
\normalsize
\renewcommand{\arraystretch}{1.2}
\adjustbox{max width=\textwidth}{%
\begin{tabular}{@{}lccc@{}}
\toprule
\textbf{$\lambda_{\text{KL}}$} & \textbf{Reward Margin} & \textbf{ECLIPTICA Score} & \textbf{Stability} \\
\midrule
0 (no KL) & 3.8 & 0.22 & Unstable \\ \hline
0.00005 & 5.1 & 0.28 & Marginal \\ \hline
0.0001 & 6.2 & 0.33 & Stable \\ \hline
0.00023 (best) & 7.5 & 0.37 & Stable \\ \hline
0.0005 & 6.8 & 0.34 & Stable \\ \hline
0.001 & 5.5 & 0.29 & Stable \\
\bottomrule
\end{tabular}%
}
\caption{\textbf{Effect of KL weight on stability and switchability.} Optimal performance occurs in a narrow band $\lambda_{\text{KL}}\approx 2\times 10^{-4}$--$3\times 10^{-4}$, where preference separation is strong while instruction regimes remain co-located.}
\label{tab:kl_ablation}
\end{table*}

\vspace{-0.35em}
\noindent\textbf{Reading Table~\ref{tab:kl_ablation}.}
\textbf{(i) No anchor collapses switching.}
At $\lambda_{\text{KL}}=0$, the model can increase margins by drifting far from the reference, which yields \textbf{highly variable training dynamics} and the lowest ECLIPTICA score (0.22).
Empirically, this manifests as \textbf{regime interference}: one instruction posture dominates, and counterfactual instruction flips produce weaker policy changes.

\vspace{-0.25em}
\noindent
\textbf{(ii) A ``Goldilocks'' KL band maximizes switchability.}
As $\lambda_{\text{KL}}$ increases from $5\times 10^{-5}$ to $2.3\times 10^{-4}$, both reward margin and ECLIPTICA score improve monotonically, peaking at (7.5, 0.37).
This is the regime where the anchor is strong enough to prevent \textbf{runaway drift}, yet weak enough to permit \textbf{instruction-conditional separation}.

\vspace{-0.25em}
\noindent
\textbf{(iii) Over-anchoring suppresses learning.}
Beyond the optimum, larger $\lambda_{\text{KL}}$ reduces reward margin and ECLIPTICA score (e.g., 0.001 $\rightarrow$ 0.29), consistent with an over-tight trust region that prevents the model from carving distinct instruction-conditioned behaviors.

\vspace{-0.15em}
\paragraph{\textbf{Finding.}}
\textbf{Too low} $\lambda_{\text{KL}}$ yields \textbf{regime collapse / interference}; \textbf{too high} $\lambda_{\text{KL}}$ yields \textbf{over-constraint}.
Switchability is maximized in a narrow band where the model learns \textbf{separation without drift}.

% -----------------------------------------------------------------------------
\subsection{Effect of Temperature ($\beta$)}
% -----------------------------------------------------------------------------

\vspace{-0.25em}
\noindent
\textbf{Why this ablation matters.}
The temperature $\beta$ controls the sharpness of the logistic contrast in the preference term.
Larger $\beta$ amplifies gradients near ambiguous pairs and increases reward margin, but it can also create \textbf{over-confident separation} that makes the model less responsive to instruction flips---i.e., it learns a \textbf{hard-coded preference wall} rather than a \textbf{switchable contract}.
Smaller $\beta$ produces smoother updates but may under-separate preferences and weaken instruction-conditioned discrimination.

\begin{table*}[ht!]
\centering
\normalsize
\renewcommand{\arraystretch}{1.2}
\adjustbox{max width=\textwidth}{%
\begin{tabular}{@{}lccc@{}}
\toprule
\textbf{$\beta$} & \textbf{Pref.\ Accuracy} & \textbf{Reward Margin} & \textbf{ECLIPTICA Score} \\
\midrule
0.05 & 85\% & 4.2 & 0.29 \\ \hline
0.10 (best) & 89\% & 7.5 & 0.37 \\ \hline
0.15 & 91\% & 8.1 & 0.35 \\ \hline
0.20 & 92\% & 9.2 & 0.31 \\
\bottomrule
\end{tabular}%
}
\caption{\textbf{Effect of temperature on preference learning and switchability.} Larger $\beta$ increases separation (margin, accuracy) but can reduce instruction-conditioned switching, revealing a margin--switchability trade-off.}
\label{tab:beta_ablation}
\end{table*}

\vspace{-0.35em}
\noindent\textbf{Reading Table~\ref{tab:beta_ablation}.}
\textbf{(i) Separation increases with $\beta$.}
Preference accuracy and reward margin grow steadily from $\beta=0.05$ to 0.20, as expected from sharper contrastive learning.

\vspace{-0.25em}
\noindent
\textbf{(ii) Switching peaks at moderate sharpness.}
ECLIPTICA score peaks at $\beta=0.10$ (0.37) and declines at larger $\beta$ despite higher margins.
This is the critical signal: \textbf{maximizing preference separation is not equivalent to maximizing instruction-conditioned control}.
Overly sharp objectives can encourage \textbf{rigid} preference boundaries that are less modulated by the instruction channel.

\vspace{-0.15em}
\paragraph{\textbf{Finding.}}
\textbf{Higher $\beta$} strengthens margins but can reduce \textbf{instruction sensitivity}.
The optimum occurs where preference learning is sharp enough to separate behaviors but smooth enough to preserve \textbf{counterfactual controllability}.

% -----------------------------------------------------------------------------
\subsection{Hyperparameter Sensitivity (Optuna/TPE)}
% -----------------------------------------------------------------------------

\vspace{-0.25em}
\noindent
We run 13 Optuna trials with TPE sampling for \textsc{CITA\_Instruct}.
TPE proposes configurations by fitting densities over \textbf{good} vs.\ \textbf{bad} trials and sampling where expected improvement is highest.
We optimize learning rate, weight decay, $\beta$, and $\lambda_{\text{KL}}$.
Figure~\ref{fig:hp_ablation_combined} summarizes the sensitivity of reward margin to these hyperparameters, revealing a concentrated ``good'' region rather than a broad plateau.

\begin{figure*}[p]
\centering
\includegraphics[width=0.85\textwidth]{figures/appendix/hp_ablation_combined.pdf}
\caption{\textbf{Hyperparameter sensitivity across 13 Optuna trials.} Reward margin is most sensitive to learning rate and the $(\beta,\lambda_{\text{KL}})$ interaction. Trial 7 ($\star$) achieves the best trade-off, aligning with the ``Goldilocks'' band: $\beta\in[0.10,0.12]$, $\lambda_{\text{KL}}\approx 2\times 10^{-4}$, and LR $\approx 5\times 10^{-6}$.}
\label{fig:hp_ablation_combined}
\end{figure*}

\paragraph{Reading the 4-panel ablation.}
Figure~\ref{fig:hp_ablation_combined} presents reward margin (y-axis) against each of the four tuned hyperparameters (x-axis): $\beta$ (DPO temperature), $\lambda_{\text{KL}}$ (KL regularization coefficient), learning rate, and weight decay.
Each blue point represents one of 13 Optuna trials, with Trial 7 ($\star$) highlighted as the best configuration.
Dashed trend lines (polynomial fits) reveal the functional relationship between each hyperparameter and reward margin.
The key observation is that reward margin is \textbf{not uniformly sensitive} to all hyperparameters: learning rate shows a strong positive correlation (higher LR $\rightarrow$ higher margin up to a point), while weight decay shows weaker sensitivity with high variance.

\paragraph{Identifying the ``Goldilocks zone''.}
Across all four panels, Trial 7 consistently appears near the peak of the trend curve, confirming that the selected configuration occupies the \textbf{intersection of optimal regions} across hyperparameters.
The $\beta$ panel shows a peaked distribution with optimum around 0.10--0.12; the $\lambda_{\text{KL}}$ panel shows an inverted-U with optimum around $2$--$3 \times 10^{-4}$; the learning rate panel shows monotonic increase up to $\sim 5 \times 10^{-6}$; and the weight decay panel shows weak sensitivity.
This multi-dimensional view explains why naive grid search would be inefficient: the optimal region is a narrow \textbf{tube} in 4D space, not a hyperrectangle.

\paragraph{Why TPE sampling works.}
The Tree-structured Parzen Estimator (TPE) efficiently navigates this complex landscape by modeling $P(\text{hyperparameters} | \text{good})$ vs.\ $P(\text{hyperparameters} | \text{bad})$.
Trials 0--4 explore broadly; trials 5--12 concentrate near promising regions.
The final distribution of points shows TPE's adaptive behavior: dense sampling near $\beta \approx 0.11$, $\lambda_{\text{KL}} \approx 2 \times 10^{-4}$, and LR $\approx 5 \times 10^{-6}$, with sparse exploration of poor regions (low LR, extreme $\beta$).

\vspace{-0.25em}
\noindent
\textbf{Pareto structure.}
Because reward margin and preference accuracy do not perfectly correlate with ECLIPTICA switchability, we additionally analyze margin--accuracy trade-offs.
Figure~\ref{fig:hp_pareto} shows the Pareto frontier; the selected configuration lies on the frontier, indicating that further margin gains typically require accuracy/switchability trade-offs.

\begin{figure}[ht!]
\centering
\includegraphics[width=0.9\columnwidth]{figures/appendix/hp_pareto_frontier.pdf}
\caption{\textbf{Pareto frontier across trials.} Trial 7 lies on the frontier, achieving near-optimal balance (margin=7.52, accuracy=89\%) while maintaining strong instruction-conditioned switching.}
\label{fig:hp_pareto}
\end{figure}

\paragraph{Interpreting the Pareto frontier.}
Figure~\ref{fig:hp_pareto} plots the two primary training objectives---reward margin (x-axis, higher is better) and preference accuracy (y-axis, higher is better)---for all 13 Optuna trials.
The red dashed line traces the \textbf{Pareto frontier}: configurations where no other trial dominates on both objectives.
Trials below or left of this frontier are Pareto-dominated (another trial is better on at least one axis without being worse on the other).
The frontier reveals a fundamental trade-off: achieving margin $>$9 (trials 0, 12) requires sacrificing accuracy to $\sim$85--87\%, while maintaining accuracy $>$89\% (trials 2, 8) limits margin to $\sim$7.

\paragraph{Why Trial 7 is optimal for switching.}
Trial 7 ($\star$) achieves margin 7.52 with accuracy 89\%, placing it \textbf{on the Pareto frontier} at the ``knee'' where both objectives are near-optimal.
Crucially, downstream evaluation shows that Trial 7 produces the best ECLIPTICA switching scores, suggesting that the margin-accuracy balance at this point corresponds to \textbf{optimal instruction-conditioned behavior}.
Trials 0 and 12 achieve higher margins (9.9, 10.1) but their lower accuracy (87\%, 86\%) correlates with degraded switching---the model over-specializes to one preference direction, losing the ability to traverse between instruction-conditioned regimes.

\paragraph{Multi-objective optimization justification.}
This Pareto analysis justifies our multi-objective Optuna formulation (maximize margin, maximize accuracy, minimize eval loss).
Single-objective optimization on margin alone would select trials 0 or 12, which underperform on switching despite high margins.
By optimizing for the Pareto frontier and selecting configurations at the knee, we ensure that \textsc{CITA} learns \textbf{balanced, traversable} preference geometry rather than collapsing to a single dominant regime.

\vspace{-0.35em}
\noindent\textbf{Trial 7 analysis (why it wins).}
The best configuration ($\beta$=0.107, $\lambda_{\text{KL}}$=0.00023, LR=5.4e-6) occupies a narrow region where:
\textbf{(i) $\beta$ is high enough} to separate instruction-conditioned preferences,
\textbf{(ii) $\lambda_{\text{KL}}$ is strong enough} to prevent regime drift,
and \textbf{(iii) LR is low enough} to avoid overshooting under longer instruction-augmented sequences.
This ``alignment'' of the three levers yields high margins without sacrificing switchability.

% -----------------------------------------------------------------------------
\subsubsection{Individual Hyperparameter Analysis}
% -----------------------------------------------------------------------------

\vspace{-0.25em}
\noindent
Figures~\ref{fig:hp_beta}--\ref{fig:hp_lr} provide per-parameter sensitivity analyses, plotting reward margin, accuracy, and evaluation loss.
Across plots, Trial 7 ($\star$) consistently appears near the optimal region, indicating that the final choice is not an outlier but a coherent solution across axes.

\begin{figure*}[p]
\centering
\includegraphics[width=0.85\textwidth]{figures/appendix/hp_ablation_beta.pdf}
\caption{\textbf{$\beta$ sensitivity.} Reward margin peaks at $\beta\approx 0.10$--$0.12$; accuracy is stable (88--90\%), while evaluation loss increases at higher $\beta$, consistent with over-sharpening. Trial 7 ($\star$) at $\beta$=0.107 achieves strong margins while preserving switchability.}
\label{fig:hp_beta}
\end{figure*}

\paragraph{Understanding $\beta$ (DPO temperature).}
Figure~\ref{fig:hp_beta} shows how the DPO temperature parameter $\beta$ affects three training metrics across 13 Optuna trials.
$\beta$ controls the \textbf{sharpness of preference learning}: higher $\beta$ increases the penalty for ranking violations, pushing the model to more decisively separate chosen from rejected responses.
The top-left panel (Reward Margin) reveals a clear \textbf{inverted-U relationship}: margin is low at extreme $\beta$ values and peaks in the 0.10--0.12 range.
At $\beta < 0.10$, preference learning is too soft and the model fails to discriminate strongly; at $\beta > 0.14$, over-sharpening causes the model to collapse toward a single dominant response pattern, destroying multi-regime switching.

\paragraph{$\beta$'s differential effects on metrics.}
The top-right panel (Accuracy) shows remarkable \textbf{stability} across the $\beta$ range: accuracy varies only from 82\% to 89\%, with no clear trend.
This indicates that $\beta$ primarily affects \textbf{separation magnitude} rather than \textbf{ranking correctness}---the model can correctly rank preferences at any $\beta$, but the confidence of that ranking varies.
The bottom panel (Eval Loss) shows an upward trend: higher $\beta$ increases loss, consistent with the over-sharpening hypothesis.
Trials 0 and 12 achieve very high margins (9.9, 10.1) but also high loss (0.51, 0.61), indicating unstable training.
Trial 7 at $\beta$=0.107 achieves strong margin (7.5) with low loss (0.33), representing the optimal balance.

\begin{figure*}[p]
\centering
\includegraphics[width=0.85\textwidth]{figures/appendix/hp_ablation_lambda_kl.pdf}
\caption{\textbf{$\lambda_{\text{KL}}$ sensitivity.} Reward margin exhibits an inverted-U: too low produces instability and regime interference; too high suppresses learning. Accuracy is comparatively stable, indicating the dominant effect is on switching stability rather than raw preference fit. Trial 7 ($\star$) at $\lambda_{\text{KL}}$=0.00023 lies in the optimal band.}
\label{fig:hp_lambda}
\end{figure*}

\paragraph{Understanding $\lambda_{\text{KL}}$ (KL regularization).}
Figure~\ref{fig:hp_lambda} shows how the KL regularization coefficient $\lambda_{\text{KL}}$ affects training metrics.
$\lambda_{\text{KL}}$ controls the \textbf{stability anchor} in \textsc{CITA}'s unified loss: it penalizes deviation from the reference policy (frozen DPO checkpoint), preventing the model from drifting too far during instruction-conditioned training.
The top-left panel (Reward Margin) shows a pronounced \textbf{inverted-U relationship}: margin peaks around $\lambda_{\text{KL}} \approx 2$--$4 \times 10^{-4}$ and drops sharply at both extremes.
At $\lambda_{\text{KL}} < 1.5 \times 10^{-4}$, the anchor is too weak to prevent regime interference, causing unstable training and low margins; at $\lambda_{\text{KL}} > 5 \times 10^{-4}$, the anchor is too strong and suppresses preference learning entirely.

\paragraph{$\lambda_{\text{KL}}$'s effect on accuracy and loss.}
The top-right panel (Accuracy) shows a \textbf{downward trend}: higher $\lambda_{\text{KL}}$ reduces accuracy from $\sim$89\% to $\sim$82\%.
This is intuitive: stronger KL regularization constrains the policy closer to the reference, limiting its ability to learn new preference patterns and thus reducing ranking accuracy.
The bottom panel (Eval Loss) shows an \textbf{upward trend}: higher $\lambda_{\text{KL}}$ increases loss, reflecting the added KL penalty term in the unified loss.
Trial 7 at $\lambda_{\text{KL}}$=0.00023 achieves the optimal trade-off: strong enough regularization to maintain stability (low loss 0.33) while preserving learning capacity (high accuracy 89\%, margin 7.5).

\paragraph{Why the KL anchor is mandatory for \textsc{CITA}.}
The sharp margin drop at low $\lambda_{\text{KL}}$ (trials 1, 3, 4) demonstrates that the KL anchor is \textbf{not optional}.
Without sufficient regularization, \textsc{CITA} training becomes unstable because the model can drift arbitrarily from the DPO-pretrained policy, losing the preference structure that enables multi-regime switching.
This finding validates a core design choice: \textsc{CITA} requires explicit anchoring to prevent the instruction-conditioning optimization from destroying the base preference geometry.

\begin{figure*}[ht!]
\centering
\includegraphics[width=0.9\textwidth]{figures/appendix/hp_ablation_learning_rate.pdf}
\caption{\textbf{Learning-rate sensitivity.} Reward margin is highly sensitive to LR, with an optimum near $5\times 10^{-6}$. Higher LR reduces accuracy and increases loss, especially for instruction-augmented sequences. \textbf{Key insight:} \textsc{CITA\_Instruct} requires $\sim$50\% lower LR than standard DPO due to longer contexts and higher gradient variance.}
\label{fig:hp_lr}
\end{figure*}

\paragraph{Learning rate is the dominant sensitivity axis.}
Figure~\ref{fig:hp_lr} reveals that learning rate has the \textbf{strongest effect} on reward margin of any hyperparameter.
The top-left panel shows a clear \textbf{monotonic positive relationship}: margin increases from $\sim$3.5 at LR=$2.5 \times 10^{-6}$ to $\sim$10 at LR=$5.4 \times 10^{-6}$.
This strong correlation (dashed trend line has high $R^2$) indicates that LR directly controls the ``aggressiveness'' of preference learning.
However, the highest-margin trials (0, 12 at LR$\approx$5.2--5.4$\times 10^{-6}$) also exhibit degraded switching on ECLIPTICA, suggesting that very high LR causes over-fitting to preference pairs at the expense of multi-regime flexibility.

\paragraph{LR's effect on accuracy and loss.}
The top-right panel (Accuracy) shows a \textbf{positive correlation}: higher LR improves accuracy from $\sim$82\% to $\sim$89\%.
Unlike $\beta$ and $\lambda_{\text{KL}}$ (where accuracy was relatively stable), LR directly affects how well the model learns to rank preferences correctly.
The bottom panel (Eval Loss) shows a \textbf{U-shaped relationship}: loss is minimized around LR=$4.5$--$5.2 \times 10^{-6}$, increasing at both lower LR (underfitting) and higher LR (overfitting/instability).
Trial 7 at LR=$5.4 \times 10^{-6}$ sits at the edge of the optimal region---slightly higher would risk instability.

\paragraph{Why \textsc{CITA\_Instruct} requires lower LR than standard DPO.}
A key insight from this ablation is that instruction-augmented training requires $\sim$50\% lower LR than NoInstruct variants.
The reason is sequence length: Instruct samples include system prompts, adding 30--40\% more tokens per sequence.
Longer sequences produce larger gradient magnitudes (more terms in the loss sum), which can cause overshooting if LR is not reduced proportionally.
Our Optuna search space accounts for this by using a conditional upper bound: \texttt{if variant == "Instruct": lr\_max = 0.5 $\times$ lr\_max\_NoInstruct}.
This constraint ensures that Trial 7's LR of $5.4 \times 10^{-6}$ is appropriate for the longer Instruct sequences.

\vspace{-0.25em}
\noindent
\textbf{Overall takeaway.}
Switchability emerges only when \textbf{separation} (controlled by $\beta$) and \textbf{stability} (controlled by $\lambda_{\text{KL}}$ and LR) are jointly satisfied.
The ablations show that maximizing reward margin alone is insufficient; the best instruction-driven alignment requires the \textbf{Goldilocks zone} where regimes remain \textbf{distinct yet co-located}, enabling reliable counterfactual switching under a fixed user prompt.





%% =============================================================================
%% EVALUATION BENEFITS
%% =============================================================================
% =============================================================================
% Appendix: Detailed Evaluation of CITA (NeurIPS spotlight; expanded)
% =============================================================================

\section{Detailed Evaluation of \textsc{CITA}}
\label{sec:appendix_benefits}

\vspace{-0.25em}
\noindent
This appendix expands the evaluation beyond aggregate scores, clarifying \textbf{what kind of capability} \textsc{CITA} adds and \textbf{how to interpret} improvements in the ECLIPTICA setting.
Our thesis is that instruction-driven alignment should be assessed as \textbf{counterfactual policy control}:
holding the \textbf{user request} fixed while changing only the \textbf{alignment instruction contract}.
Accordingly, we organize the evaluation into four operational properties that matter in deployment:
\textbf{generalization}, \textbf{robustness}, \textbf{over-refusal calibration}, and \textbf{policy switching}.
For each, we state (i) what we measure, (ii) what success looks like, and (iii) what failure modes look like.

\vspace{-0.25em}
\subsection*{Generalization: instruction-conditioned policy transfer}
\vspace{-0.15em}

\noindent
\textbf{Claim.} \textsc{CITA} learns an instruction-conditioned policy family that \textbf{transfers} to prompts and topics not observed during training.

\vspace{-0.15em}
\noindent
\textbf{What we measure.}
We evaluate whether an alignment instruction $I$ induces the intended behavioral posture on \textbf{novel prompts}:
(i) unseen domains (e.g., finance $\rightarrow$ health advice posture, crypto $\rightarrow$ workplace communication),
(ii) different surface forms (question, imperative, multi-part),
(iii) distributional shifts in topic and framing.

\vspace{-0.15em}
\noindent
\textbf{What success looks like.}
A generalized policy does \textbf{not} memorize prompt templates.
Instead, it applies \textbf{contract-consistent content selection}:
a \texttt{safety\_first} instruction should reliably increase warnings and guardrails,
a \texttt{concise} instruction should reliably compress while preserving key decision variables,
and a \texttt{professional} instruction should produce consistent structure and hedging behavior.
Importantly, this should hold even when the prompt is phrased adversarially or indirectly.

\vspace{-0.15em}
\noindent
\textbf{What failure looks like.}
We observe three typical generalization failures in instruction-conditioned systems:
\textbf{(i) template lock-in} (works on familiar formats, fails on paraphrases),
\textbf{(ii) topic leakage} (instruction affects only some domains),
\textbf{(iii) shallow style compliance} (tone changes but policy does not).
ECLIPTICA’s fixed-prompt counterfactual design is explicitly meant to detect (iii).

\vspace{-0.25em}
\subsection*{Robustness: resisting adversarial instruction and prompt attacks}
\vspace{-0.15em}

\noindent
\textbf{Claim.} A switchable alignment mechanism must remain stable under \textbf{jailbreak-like perturbations} and \textbf{adversarial rephrasings}.

\vspace{-0.15em}
\noindent
\textbf{Threat model (high-level).}
We consider attacks that attempt to:
\textbf{(i) override} the alignment instruction (e.g., ``ignore previous policy''-style attempts),
\textbf{(ii) blur} the user intent to bypass refusal boundaries,
\textbf{(iii) induce} instruction collisions (conflicting constraints),
or \textbf{(iv) smuggle} policy changes via the user prompt.

\vspace{-0.15em}
\noindent
\textbf{What we measure.}
Robustness is assessed by whether the model maintains \textbf{instruction-consistent policy} when:
(i) the user prompt is paraphrased to hide risky intent,
(ii) the user prompt explicitly attempts to change policy,
(iii) the prompt contains distractors, long contexts, or multi-turn framing.

\vspace{-0.15em}
\noindent
\textbf{What success looks like.}
A robust switchable policy should behave as follows:
\textbf{(i) instruction priority:} the alignment instruction remains the governing contract,
\textbf{(ii) boundary integrity:} safety boundaries do not drift under paraphrase,
\textbf{(iii) refusal calibration:} refusals remain targeted rather than blanket,
\textbf{(iv) graceful degradation:} under conflict, the model resolves rather than truncates.

\vspace{-0.15em}
\noindent
\textbf{What failure looks like.}
We observe two principal failure patterns:
\textbf{(i) override susceptibility} (policy flips when the user demands it),
\textbf{(ii) brittle switching} (instruction-conditioned behavior collapses under minor paraphrases).
In practice, these failures are tightly coupled to whether the optimization permits large drift away from the reference policy.

\vspace{-0.25em}
\subsection*{Over-refusal calibration: avoiding unnecessary refusals}
\vspace{-0.15em}

\noindent
\textbf{Claim.} Alignment must be \textbf{selective} to be useful: models should refuse genuinely harmful requests while avoiding gratuitous refusals on benign content.

\vspace{-0.15em}
\noindent
\textbf{What we measure.}
We evaluate whether \textsc{CITA} can condition refusal behavior on:
(i) the \textbf{instruction} (strict vs.\ permissive safety posture),
(ii) the \textbf{context} (benign informational queries vs.\ actionable harm),
(iii) the \textbf{degree of risk} (low-risk educational vs.\ high-risk operational).

\vspace{-0.15em}
\noindent
\textbf{What success looks like.}
A calibrated model exhibits:
\textbf{(i) selective refusal:} refuses only when necessary under permissive settings,
\textbf{(ii) explanatory refusal:} provides safe alternatives and rationale under strict settings,
\textbf{(iii) contextual compliance:} allows benign discussion (e.g., safety education) while refusing harmful operational detail.

\vspace{-0.15em}
\noindent
\textbf{What failure looks like.}
Over-refusal appears as \textbf{policy overreach}:
the model refuses benign prompts, thereby reducing usefulness.
Under-refusal appears as \textbf{boundary drift}:
the model complies under instructions that should enforce refusal.
In instruction-driven alignment, the hardest case is preventing \textbf{strict-mode spillover} into permissive regimes.

\vspace{-0.25em}
\subsection*{Policy switching: counterfactual control under a shared backbone}
\vspace{-0.15em}

\noindent
\textbf{Claim.} \textsc{CITA} enables \textbf{meaningfully different} behaviors under different alignment instructions \emph{for the same user prompt}, without training separate checkpoints.

\vspace{-0.15em}
\noindent
\textbf{What we measure.}
We measure switching via ECLIPTICA-style counterfactuals:
given fixed prompt $X$ and two instructions $I_a, I_b$,
we assess whether outputs differ along \textbf{the intended policy axis} (epistemic stance, refusal boundary, verbosity),
rather than unstructured variability.
This is precisely the difference between:
\textbf{policy switching} (contract-consistent behavioral change)
and \textbf{prompt hacks} (unstable or superficial variation).

\vspace{-0.15em}
\noindent
\textbf{What success looks like.}
A strong switch is:
\textbf{(i) directional:} changes align with the instruction contract,
\textbf{(ii) consistent:} repeats across paraphrases and domains,
\textbf{(iii) bounded:} remains within a stable neighborhood of the reference (no runaway drift),
\textbf{(iv) non-degenerate:} does not collapse into trivial acknowledgements or truncated replies.

\vspace{-0.15em}
\noindent
\textbf{What failure looks like.}
We observe three failure types:
\textbf{(i) mode collapse} (one regime dominates regardless of instruction),
\textbf{(ii) regime interference} (switching works but bleeds constraints across regimes),
\textbf{(iii) brittle conjunction} (multi-instruction composition produces incomplete outputs).
The failure case in Appendix~\ref{sec:appendix_examples} illustrates (iii).

\vspace{-0.15em}
\paragraph{\textbf{Practical summary.}}
Across these dimensions, \textsc{CITA} is best understood as providing a \textbf{deployable control interface}:
alignment becomes a \textbf{runtime instruction channel} that modulates policy with bounded, repeatable counterfactual effects,
rather than a one-shot, training-time imprint of a single behavior.

% =============================================================================
% TRAINING CURVES (SUPPLEMENTARY)
% =============================================================================
\section{Training Curves}
\label{sec:appendix_training_curves}

\vspace{-0.25em}
\noindent
This section reports supplementary training dynamics that contextualize the main results in Section~\ref{sec:training_dynamics}.
We emphasize one operational point: \textbf{loss scale is not comparable across methods}.
Offline preference objectives (DPO/\textsc{CITA}) and online RL objectives (GRPO) produce different loss magnitudes and noise characteristics.
Accordingly, we interpret curves through \textbf{convergence shape}, \textbf{stability}, and \textbf{relative trends} (Instruct vs.\ NoInstruct), rather than absolute loss values.

\subsection{DPO/\textsc{CITA} Eval Loss}

\begin{figure*}[ht!]
\centering
\includegraphics[width=0.9\textwidth]{figures/training/combined_eval_loss.pdf}
\caption{\textbf{Evaluation loss across policy optimization methods.}
We plot evaluation loss for DPO, \textsc{CITA}, and GRPO using independent y-axes due to non-comparable loss scales.
\textbf{DPO (left):} loss $\sim$0.21--0.28 with smooth convergence; the Instruct variant is slightly lower, consistent with stronger preference fit under the instruction-augmented setup.
\textbf{\textsc{CITA} (center):} loss $\sim$0.27--0.40; the NoInstruct curve converges lower, while \textsc{CITA\_Instruct} attains larger reward margins, reflecting a separation--stability trade-off induced by the mandatory anchoring term.
\textbf{GRPO (right):} loss oscillates near zero ($\pm$0.001), characteristic of online RL training where the objective is dominated by on-policy sampling variance.
\textbf{Note:} PPO tensorboard logs were empty due to a TRL logging issue, so we do not include PPO curves here.}
\label{fig:training_loss}
\vspace{-3mm}
\end{figure*}

\vspace{-0.35em}
\noindent\textbf{Key takeaways from Figure~\ref{fig:training_loss}.}
\textbf{First,} DPO exhibits the most stable and smooth offline convergence, serving as a strong baseline for preference fitting.
\textbf{Second,} \textsc{CITA} exhibits slightly higher eval loss but achieves higher reward margins, consistent with learning a \textbf{switchable policy family} rather than optimizing a single preference boundary.
\textbf{Third,} GRPO's oscillatory loss highlights the cost of online methods: additional variance and sensitivity to sampling, even when downstream performance can be competitive.

\paragraph{DPO loss dynamics (left panel).}
DPO's evaluation loss decreases monotonically from $\sim$0.28 to $\sim$0.21 over 1,400 training steps, exhibiting textbook offline optimization behavior.
The Instruct variant achieves slightly lower final loss (0.21) compared to NoInstruct (0.22), indicating that instruction-augmented training provides marginally stronger preference signal.
This smooth convergence reflects DPO's core design: by reformulating RLHF as supervised learning on preference pairs, DPO avoids the sampling variance inherent to online methods.
The monotonic decrease also suggests that the learning rate ($1 \times 10^{-5}$) and LoRA configuration (r=16) are well-calibrated for this task.

\paragraph{\textsc{CITA} loss dynamics (center panel).}
\textsc{CITA} exhibits a different pattern: loss decreases from $\sim$0.40 to $\sim$0.27 (NoInstruct) or $\sim$0.34 (Instruct).
The higher absolute loss compared to DPO is expected because \textsc{CITA}'s unified loss includes both the DPO term and the KL anchor: $\mathcal{L}_{\text{CITA}} = \mathcal{L}_{\text{DPO}} + \lambda_{\text{KL}} \cdot \mathcal{L}_{\text{KL}}$.
The KL term penalizes deviation from the reference policy, which necessarily increases total loss while providing stability benefits.
Notably, \textsc{CITA}\_NoInstruct converges to lower loss than \textsc{CITA}\_Instruct---the opposite of DPO's pattern---because longer Instruct sequences incur proportionally larger KL penalties.

\paragraph{GRPO loss dynamics (right panel).}
GRPO's loss oscillates near zero ($\pm$0.001), which is characteristic of online RL rather than a sign of instability.
Unlike offline methods that compute loss on a fixed preference dataset, GRPO generates fresh responses at each step and computes group-relative advantages within each batch.
This on-policy sampling introduces inherent variance: different batches of generated responses produce different advantage distributions.
The oscillation around zero reflects that GRPO's loss is a policy gradient objective (expected advantage-weighted log-probability) rather than a supervised loss, making direct magnitude comparisons with DPO/\textsc{CITA} inappropriate.
GRPO also runs for fewer steps ($\sim$500) due to the computational overhead of online generation.

\paragraph{Why PPO curves are absent.}
PPO tensorboard logs were empty due to a known TRL logging issue where the PPOTrainer's internal metrics are not properly exported to tensorboard in certain configurations.
Despite the missing training curves, PPO evaluation results (Section~\ref{sec:overall_results}) confirm that PPO training completed successfully and produced competitive models, particularly on the AQI benchmark where PPO's online optimization shapes global reward-aligned behavior.


